% Section 7: Wood points as a limitation, with the fallback used by the search.
% Sources: test/bie-bound/wood-probe-v1 (mechanism), failure-probe-v1 P1
% (observed failures), wood-fallback-v1 (fallback and delta rule).
\section{Wood points}\label{sec:wood}
At a Wood point $k_W$ some retained order has $\gamma_m(k_W)=0$. The
guided-mode problem is regular there (\cref{sec:rayleigh}), but the solver of
\cref{sec:discretization,sec:leads,sec:closure} is not. This section
explains where it fails, and describes how the search avoids these points.
We do not remove the singularity in this paper.

\subsection{Where the solver fails}\label{sec:wood-mechanism}
Let $t=|k-k_W|$ and take the order $m=0$ with $\rho_{\mathrm e}=1$, so that
$|\gamma_0|\approx(2k_Wt)^{1/2}$. The measurements below are for the periodicity
and defect cells of Example 1 at $\beta=1.98154$, for $10^{-7}\le t\le10^{-1}$ on
both sides of $k_W$ and on two grids (\texttt{wood-probe-v1}).
\begin{enumerate}
\item \emph{The cell equation.} The $m=0$ term of the spectral series
  \eqref{eq:gqp-spectral} carries the factor $1/\gamma_0$, so the
  quasi-periodic Green function itself diverges at $k_W$. The fitted exponent
  of this term against $t$ is between $-0.5025$ and $-0.5000$. The norm of
  the Müller matrix \eqref{eq:cell-system} grows with exponent about $-0.50$,
  while its smallest singular value stays constant (exponents in
  $[-0.003,0.001]$).
\item \emph{Extraction.} The extraction kernels \eqref{eq:extract-kernels}
  divide by $\gamma_m$. At $t=0$ this is a division by zero, which is where
  the code first fails.
\item \emph{The cell scattering matrix.} $S_{\mathrm{cell}}(k)-S_{\mathrm{cell}}(k_W)$
  behaves like $t^{1/2}$ on each side; the fitted exponents of the matrix
  norm are $0.496$ to $0.500$. This is the square-root branch of the Wood
  anomaly of an open grating \cite{HesselOliner1965,WojcikWangOrensteinFan2021}.
  Interpolating the determinant in $k$ across $k_W$ therefore converges only
  algebraically.
\end{enumerate}
The third item is a property of the amplitude coordinates. Convert the cell
map to wall values and normal derivatives, the coordinates
\eqref{eq:cauchy}, and form the two-wall Dirichlet-to-Neumann map of the cell.
It is linear in $k$ across $k_W$ (fitted exponents $1.000$ to $1.003$), and
the bordered determinant built from it has Chebyshev coefficients that
decrease by a factor of about $0.011$ per degree. The conversion does not
repair the first item, however. The Green function evaluated through the
proxy fit differs from the spectral series by about $1.4\times10^{-14}/t$
($1.36\times10^{-7}$ at $t=10^{-7}$, $1.2\times10^{-10}$ at $t=10^{-4}$), and the
Chebyshev coefficients of the converted determinant level off at about
$4\times10^{-10}$. This error is the same on both density grids, because they
share the Green precomputation, so the two-grid comparison does not detect
it.

A solver that is regular at $k_W$ must therefore change the cell problem, not
only its output coordinates: the cell representation must avoid the
divergent $m=0$ term, and the closure must use coordinates such as
\eqref{eq:cauchy} in which the threshold solutions $1$ and $x-x_p$ are both
present. Integral equations of this kind exist for open gratings
\cite{BarnettGreengard2011,ChoBarnett2015,BrunoDelourme2014,PerezArancibiaShipmanTurcVenakides2019}
and for band structures of doubly periodic crystals
\cite{BarnettGreengard2010,PerezArancibiaYin2026}. Combining one with the Floquet
eigenvalue problem and the closure of this paper is left to future work.

\subsection{Fallback in the search}\label{sec:wood-fallback}
Wood points are known before any solve: $k_W=|\beta_m|/\rho_{\mathrm e}^{1/2}$ for
the retained orders. The search treats them as follows.
\begin{enumerate}
\item Split the window $[k_1,k_2]$ at every $k_W$ inside it.
\item Remove the band $|k-k_W|<\delta$ from each piece, and never evaluate the
  cell matrices there.
\item On the panels next to $k_W$, interpolate the determinant in
  $u=\operatorname{sign}(k-k_W)\,|k-k_W|^{1/2}$ instead of $k$, which follows the
  square-root branch found in \cref{sec:wood-mechanism}. Away from $k_W$ the two variables give
  the same roots; near $k_W$ the variable $u$ needed fewer evaluations
  (\cref{sec:exp-wood}).
\item Report every removed band with the result. A guided mode in a removed
  band is reported as \emph{not searched}, not as absent.
\end{enumerate}
This meets requirement 5 of \cref{sec:task}, and it gives up requirement 2
inside the removed bands. Along a dispersion curve $k_*(\beta)$ a crossing of
$k_W(\beta)$ is a single point. The fallback misses $k_*$ only for the
$\beta$ in a small interval around the crossing; for these $\beta$ one can
report $k_*$ by interpolation from neighboring $\beta$, labelled as such.

Roots close to a Wood point are less accurate than others, because the
proxy error above enters the determinant, and the two-grid comparison does
not detect it. In the experiments of \cref{sec:exp-wood} a root at distance
$t$ from $k_W$ had an error of at most about $4\times10^{-14}/t$. For a target
error $\varepsilon$ we therefore choose
\[
  \delta=\max\bigl(10^{-3},\ 5\times10^{-14}/\varepsilon\bigr).
\]
The first term is the smallest width at which every tested search finished;
the second rounds up the measured constant. Both are empirical, measured on
two examples with $\rho_{\mathrm e}=1$ and $10^{-4}\le t\le10^{-1}$; the rule is not
a bound. A smaller $\delta$ may be tried, but every strip it adds must pass
the interpolation tests, and a strip that does not finish is reported as not
searched.
